\section{Optional research result: both parities of the anchor exponent}

The unchanged starting construction is the sixth-round package, whose
builder has SHA256
\texttt{3f63cb782d52857446cf39f6e545aeb74f380d1347826eafa888cbb487738d51}.
The upstream source remains pinned at
\texttt{036cebe4f1a4aa7ca89820a9b48f3629b4c364d6}.
All curve, rank, quaternion, height, pole, simultaneous-prime and Lucas
inputs retain their previously stated qualifications.  The assertions
below are new applications and optimizations of those inputs, not claims
that a historical proof lacked justification.

The actual local arguments used here are the upstream
\texttt{subsec:o8-anchor-subgroup}, lines 8132--8244 of
\texttt{manuscripts/quartic\_v10.tex}; its quaternion criterion at
lines 10973--11024; the leading-residue and product-parity argument in the
previously delivered \texttt{packed\_anchor\_note\_body.tex}; and the
sixth-round stored-root anchor proof.  The actual final height consumers
are \texttt{docs/quartic-523-540/addition.tex}, lines 421--477.

Write $K=\mathbb Q(\rho)$, $\rho^2=2$, and
$N(x+y\rho)=x^2-2y^2$.  The four retained root points are
$Q_i=[n_i]P$, one of which is the scalar root $S=[a]P$.  Put
\[
 C=\prod_i x(Q_i)=C_0+C_1\rho,\qquad B_0=4096C_0,
 \qquad Q=S-P=[a-1]P.
\]
The unchanged scalar chart and nonzero guard give
$a\ne0$, $a-1\ne0$, $h_S\ne0$, and the already stored field coordinate
\[
 w=h_S\tau,\qquad h_S^2=x(S),\qquad\tau^2=f(Q).
\]
These are exact identities on the retained domain, including the
permitted tangent case; they are not only local congruences.

Let $E\ge1$ be any integer and set
\[
 S_0=3^E,\qquad \epsilon=1-(E\bmod2),\qquad
 \gamma=3^\epsilon,\qquad
 \Theta=N(w)+\gamma S_0N(B_0+X\rho).
\]
Use precisely the same three-witness conic
\begin{equation}\label{r7:anchor}
 u^2+6v^2-2z^2-3\Theta^2+2S_0^2=0.
\end{equation}
For even $E$ the coefficient $\gamma$ is the preceding coefficient $3$;
for odd $E$ it is $1$.  No new coordinate is required.  The stored
$C_0^2$ and input-square coordinates keep the physical expression
quadratic when $S_0$ is lifted to a coordinate.

\subsection{An exact point/input-domain characterization}

Dividing the conic by $S_0^2$ and applying the source quaternion criterion
identifies existence of its rational witnesses with simultaneous
$2$- and $3$-integrality of
\[
 t=\Theta/S_0=N(w)/3^E+\gamma N(B_0+X\rho).
\]
The criterion includes $t=0$ and imposes no real sign; $S_0$ need not
be a square.

At the unramified place above $3$ the norm is anisotropic modulo $3$:
\[
 v_3(N(x+y\rho))=2\min(v_3x,v_3y).
\]
The valuations of the two displayed summands of $t$, when nonzero, have
parities $-E$ and $\epsilon$.  These parities are opposite.  Therefore
$t$ is $3$-integral exactly when
\[
 v_3(w_0),v_3(w_1)\ge\lceil E/2\rceil,
 \qquad B_0,X\in\mathbb Z_{(3)}.
\]
The second equivalence holds for both $\epsilon=0$ and $\epsilon=1$,
since $\epsilon+2m\ge0$ for an integer $m$ is equivalent to $m\ge0$.
Zero summands do not affect this argument.

The actual product leading-residue proof gives, before any root parity
assumption,
\[
 v_3(C_0)=-\sum_i e(n_i),\qquad
 e(n)=\begin{cases}0&n\text{ odd},\\2+2v_3(n)&n\text{ even}.
 \end{cases}
\]
Thus the inferred integrality of $B_0$ forces all four labels to be odd.
In particular $a$ is odd.  The reduction of $P$ at $3$ has order two,
so $S$ reduces to $P$ and $h_S$ is a unit.  The source's formal multiple
calculation now gives
\[
 v_3(w)=v_3(\tau)=1+v_3(a-1).
\]
The conclusion at $3$ is therefore exactly oddness of all four labels,
$X\in\mathbb Z_{(3)}$, and
\begin{equation}\label{r7:index}
 0\ne a-1\in2\cdot3^{\lceil E/2\rceil-1}\mathbb Z.
\end{equation}
This order of deductions avoids circular use of the odd-product height
identity.

At $2$, normalize $v(2)=1$ and $v(\rho)=1/2$.  The integral odd formal
series, the depth $v(z(P))=3/2$, and
$[2]z=2z+O(z^3)$ imply
\[
 v(z([n]P))=3/2+v_2(n),\qquad
 v(x([n]P))=-3-2v_2(n)
\]
for every nonzero integer $n$.  The higher terms are strictly deeper
because $3r>r+1$ for $r\ge3/2$; odd multiplication preserves depth.
The already obtained odd labels make $B_0$ $2$-integral and give
\[
 v(h_S)=-3/2,\qquad v(\tau)=3/2+v_2(a-1),\qquad
 v(w)=v_2(a-1)\ge1.
\]
The rational basis valuations have distinct fractional parts, so both
$w_0,w_1$ are $2$-integral.  Since $3^E$ and $\gamma$ are odd, all
terms of $t$ except $-2\gamma X^2$ are now $2$-integral.  If $X$ is
not $2$-integral, that last term has negative valuation
$1+2v_2(X)$ and cannot cancel.  The remaining conic condition is thus
exactly $X\in\mathbb Z_{(2)}$.

Conversely these same point/input conditions make $t$ integral at both
primes, and the quaternion criterion supplies its witnesses.  This
proves the uniform equivalence
\[
 \boxed{\eqref{r7:anchor}\quad\Longleftrightarrow\quad
 \text{all four labels odd},\quad
 X\in\mathbb Z_{(2)}\cap\mathbb Z_{(3)},\quad\eqref{r7:index}.}
\]
It is an equivalence of existence of conic witnesses, not a bijection
between witness tuples.  All sign choices of the retained scalar chart
are covered by its squared identities and by the norm.

In particular odd exponent $2k-1$ with $\gamma=1$ and even exponent
$2k$ with $\gamma=3$ have exactly the same point/input domain.
Completeness still permits arbitrary divisibility of the positive
scalar index $a-1$, so either exponent imposes only a finite preparation
requirement and does not alter the order of the later prime choices.

\section{Projective height of the three retained scales}

This argument sharpens a constant without changing the scales or their
source equations.  Put $U=C_0+C_1$.  The three rational scales are
$1,C_0^2,U^2$; their established bounds $C_0\ge193$ and $U\ge61$
make them positive and nonzero on every source solution.  If a nonzero
integer $N$ is divisible to order $k$ by the three contact radicals,
the existing detector and fractional-linear height estimates give
\[
 h(\lambda),\quad h(C_0^2\lambda),\quad h(U^2\lambda)
 \ \le T=(H+\log50)|N|^{19/k}.
\]
Use absolute logarithmic projective height.  Its scale invariance and
the exact height rule for coordinatewise squaring give
\[
 2h([1:C_0:U])
 =h([\lambda:C_0^2\lambda:U^2\lambda])
 \le h(\lambda)+h(C_0^2\lambda)+h(U^2\lambda)
 \le3T.
\]
The middle inequality follows place by place from
$\log\max_i|y_i|_v\le\sum_i\log\max(1,|y_i|_v)$.
All its inputs are nonzero, so no undefined projective tuple occurs.

Now $C=(1-\rho)C_0+\rho U$.  At finite places the coefficients
$1-\rho$ and $\rho$ are integral.  At the two real places the operator
bounds relative to $\max(1,|C_0|,|U|)$ are respectively
$2\sqrt2-1$ and $2\sqrt2+1$, whose product is $7$.  Hence
\[
 h(C)\le h([1:C_0:U])+\tfrac12\log7
 \le\tfrac32 T+\tfrac12\log7.
\]
The actual odd-product argument in
\texttt{docs/quartic-523-540/addition.tex}, lines 296--350, bounds each
root height by $h(C)$; this is the product in $K$, not its rational norm.
With $B=\max(|n|,|a|,|t|,|b|)$ its inherited lower bound therefore yields
\begin{equation}\label{r7:height}
 B^2-6 < K_0 |N|^{19/k},\qquad K_0=(3H+14)/2.
\end{equation}
Here $\log50<4$, $\log7<2$, and $|N|^{19/k}\ge1$ justify the strict
rational constant $K_0$.  No new height theorem is being invoked.
The old completeness proof requires only fixed positive rational-square
scales after the roots have been chosen; all three scales are unchanged.

\section{Uniform numerical consumers}

Retain the inherited finite-exception height constant
\[
 H=2^{110}(2^{32}3^9)^6\,46410^{19}+128.
\]
Its exact bit length is $683$.  For a scalar order $m>19$, anchor
exponent $E$, and $q=\lceil E/2\rceil$, define
\[
 c_*=2\cdot3^{q-1},\qquad b_*=c_*-1.
\]
Every source solution has $c=|a-1|\ge c_*$ and $B\ge|a|\ge b_*$.
The unchanged sufficient scalar margin follows from the exact integer
test
\begin{equation}\label{r7:scalar-test}
 c_*^{m-19}>(80H)^m.
\end{equation}
It gives $h(\beta)<c/80\le|a|/40$.  The separate logarithmic margin
$\log(40H)<|a|/40$ follows, for the parameters considered here, from
$H<2^{683}$ and $40\cdot689<b_*$.  This uses $40<2^6$ and
$\log2<1$.

For the literal Lucas polynomial
\[
 F=t^2-a(a-2)tb+(a-2)^2b^2-1,
\]
the source triangle bound is
$|F|\le P(B)=2B^4+6B^3+5B^2+1$.
At high order $k>38$, the contradiction $F\ne0$ is therefore certified
by the exact integer inequality
\begin{equation}\label{r7:lucas-test}
 [2(b_*^2-6)]^k>(3H+14)^k P(b_*)^{19}.
\end{equation}
This finite test implies the required inequality for every $B\ge b_*$,
not just at the endpoint: for $B>\sqrt6$,
\[
 \frac{(B^2-6)^k}{P(B)^{19}}
 =B^{2k-76}
   \frac{(1-6/B^2)^k}{(2+6/B+5/B^2+B^{-4})^{19}}
\]
is strictly increasing, since each positive factor is nondecreasing and
the first has positive exponent.  Thus \eqref{r7:height} with $N=F$
is impossible.  The inherited Lucas conclusion is
$\log B>|a|-3$.

The input consumer remains valid at any low order $\ell\ge10$ supported
by its logarithm.  Since $K_0<7H$, the old estimate remains available:
if $n\ne\beta$, its reduced numerator satisfies
$|N|\le2e^{h(\beta)}B$, and
\[
 \log B<\frac{\ell}{2\ell-19}\log(40H)
             +\frac{19}{2\ell-19}h(\beta).
\]
For $\ell=10,11,12$ the sum of the two coefficients is at most $29$,
so the last expression is less than $29|a|/40<|a|-3$ once
$11b_*>120$.  This checks the input consumer separately from the
scalar and Lucas margins.  A lower order unsupported by this height
argument is not thereby asserted impossible for every alternative
argument.

These exact tests are sufficient conditions, not optimality claims.
Larger $E$ may compensate for smaller scalar order $m>19$ or high
order $k>38$.  The companion certificate reports successful and failed
finite tests distinctly; failure means only that the displayed
sufficient inequality has not been established for those parameters.


\subsection{Illustrative sufficient exponents}

These refinements are not dependencies of the selected exponent $2790$.
With scalar order $28$, high order $45$ and input order $12$, exact
integer evaluation of the displayed sufficient tests first succeeds at
odd $E=2771$ (equivalently even $E=2772$).  The earlier $24H$ Lucas
bound first succeeds at odd $E=2787$ (equivalently even $E=2788$).
These are minima only for the specific sufficient integer comparisons
with the stated constants.  No circuit, coefficient-size, or global
mathematical optimality is asserted.  The selected even exponent $2790$
has the simpler identity $3^{2790}=(3^{465})^6$ and already satisfies the
preceding constants, so it does not rely on either optional refinement.
